\chapter{相关理论与技术基础}
\vspace{7pt}


\section{分子表示方法}

在计算机辅助药物设计中,如何把分子的化学结构用计算机能够处理的形式来表示,是一个最基础也是最关键的问题。分子的化学结构本质上是一个图结构,原子是节点,化学键是边,这种图结构包含了丰富的拓扑信息和几何信息。为了能够在计算机上存储、处理和比较分子,研究者开发了多种分子表示方法,每种方法都有其适用的场景和优缺点。

最早也是最直观的表示方法是连接表(Connection Table),它用一个表格来记录分子中所有的原子和键的信息。每一行代表一个原子,记录原子的类型、坐标、电荷等属性；另外的表格记录键的信息,包括连接的原子、键的类型(单键、双键、三键等)。这种表示方法非常详细,能够完整地描述分子的结构,但是它占用的存储空间比较大,而且不太适合直接用于机器学习的输入。

为了解决这个问题,线性记号系统应运而生,其中最有代表性的就是SMILES(Simplified Molecular Input Line Entry System)。SMILES是由David Weininger在1988年提出的\cite{weininger1988smiles},它的核心思想是用一串ASCII字符来表示分子的结构。具体来说,原子用元素符号表示,比如C代表碳原子,O代表氧原子,N代表氮原子等；键用特定的符号表示,单键可以省略,双键用\texttt{=}表示,三键用\texttt{\#}表示；环结构通过断键并添加数字标记来表示,比如c1ccccc1表示苯环；支链用括号括起来。举个例子,乙醇的SMILES是\texttt{CCO},表示两个碳原子和一个氧原子依次连接；乙酸的SMILES是\texttt{CC(=O)O},表示一个甲基连接到一个羧基上。

SMILES表示有几个明显的优点。一是简洁紧凑,一个简单的分子只需要几个到几十个字符就能表示,非常适合存储和传输；二是可读性较好,经过训练的化学家可以直接从SMILES字符串看出分子的大致结构；三是易于处理,作为字符串,它可以方便地用于文本处理、模式匹配、以及机器学习模型的输入。不过SMILES也有一些局限性,比如同一个分子可能有多种合法的SMILES表示(这被称为SMILES异构性),不同的SMILES字符串可能对应同一个分子,这会给比较和去重带来困难；另外,SMILES主要描述的是二维的拓扑结构,对于三维的空间构象信息表达不够充分。

除了SMILES,还有其他一些线性记号系统,比如InChI(International Chemical Identifier)、SELFIES(Self-Referencing Embedded Strings)等。InChI是由国际纯粹与应用化学联合会(IUPAC)制定的标准,它的特点是标准化程度高,同一个分子只会生成唯一的InChI字符串,解决了SMILES的异构性问题。不过InChI的可读性比较差,人类很难直接从InChI字符串看出分子的结构。SELFIES是近年来提出的一种新的表示方法,它的设计目标是保证生成的任何字符串都对应一个有效的分子,这对于生成模型来说非常重要,因为可以避免生成无效的分子结构。

除了线性的字符串表示,还有一种重要的表示方法是分子指纹(Molecular Fingerprint)。分子指纹是把分子结构编码成一个固定长度的二进制向量或者实数向量,向量中的每一位代表某种结构特征是否存在。常见的分子指纹有MACCS密钥、ECFP(Extended Connectivity Fingerprints)、PubChem指纹等。以ECFP为例,它是一种圆形指纹,通过迭代地考虑每个原子周围的邻居环境,生成一系列的特征标识符,然后把这些标识符哈希到一个固定长度的向量中。分子指纹的优点是可以直接用传统的机器学习算法进行处理,比如支持向量机、随机森林等,而且计算相似度非常方便,可以用Tanimoto系数等指标来衡量两个分子的相似程度。缺点是它是基于预定义的特征,表达能力有限,而且丢失了部分结构信息。

最近几年,随着深度学习的发展,出现了很多端到端的分子表示学习方法。这些方法不需要手工设计特征,而是让模型自动从数据中学习分子的表示。比如图神经网络(Graph Neural Network, GNN),它直接把分子当作图来处理,通过消息传递机制聚合邻居节点的信息,学习每个原子的嵌入表示,然后通过池化操作得到整个分子的表示。这种方法能够充分利用分子的图结构信息,在很多分子性质预测任务上都取得了很好的效果。还有基于Transformer的方法,它把SMILES当作序列来处理,利用自注意力机制学习序列中各个位置之间的依赖关系,也能够学到很好的分子表示。

在本研究中,我们主要使用SMILES作为分子的表示方法,原因有几个方面。一是SMILES是最广泛使用的分子表示之一,大多数的化学信息学工具和数据库都支持SMILES格式,便于数据的获取和处理；二是SMILES作为字符串,可以直接作为大语言模型的输入和输出,不需要额外的转换或编码；三是SMILES的语法相对简单,大语言模型比较容易学习和理解；四是我们使用了RDKit库来进行SMILES的标准化和验证,可以有效地处理SMILES异构性和有效性检查的问题。当然,SMILES的局限性也是存在的,比如在后续的改进工作中,可以考虑结合其他表示方法,或者探索更适合大语言模型处理的分子表示形式。

\section{药物属性预测}

在药物分子设计中,评估一个分子的潜在价值需要计算多个药物相关的属性指标。这些属性反映了分子在药代动力学、毒理学、合成可行性等方面的特性,是判断一个分子是否适合作为药物候选物的的重要依据。常用的药物属性包括QED(定量估计药物相似性)、LogP(辛醇-水分配系数)、分子量、氢键供体/受体数量、可旋转键数量、极性表面积等。下面我们来介绍几个关键的属性及其计算方法。

QED(Quantitative Estimate of Drug-likeness)是一个综合性的药物相似性指标,它是由Bickerton等人在2012年提出的\cite{bickerton2012qed}。QED的计算基于八个关键的分子描述符,包括分子量、脂溶性(LogP)、氢键供体数量、氢键受体数量、极性表面积、可旋转键数量、芳香环比例、以及警报子结构的存在情况。每个描述符都有一个理想的范围,偏离这个范围会降低分数。QED的值在0到1之间,越接近1表示分子的药物相似性越好。相比于简单的规则如Lipinski五规则,QED的优势在于它是一个连续的、平滑的函数,能够更细致地区分不同分子的药物潜力,而且它考虑的因素更全面,不仅仅局限于口服药物的要求。在我们的研究中,QED被选为主要的优化目标,因为它能够综合反映分子的多个方面的性质,是一个比较全面的评价指标。

LogP(Octanol-Water Partition Coefficient)是辛醇-水分配系数的对数值,它反映了分子的脂溶性。LogP的值越大,表示分子越容易溶解在脂质中,越容易穿过细胞膜；LogP的值越小,表示分子越亲水。对于口服药物来说,LogP的理想范围通常在0到5之间,过低会导致吸收不良,过高会导致代谢过快或者毒性增加。LogP的计算方法有很多,包括基于片段的方法(如XLogP、ALogP)、基于原子贡献的方法、以及基于机器学习的方法等。在RDKit中,提供了多种LogP的计算实现,我们可以根据需要进行选择。

分子量(Molecular Weight)是一个基本的物理化学属性,它直接影响分子的吸收、分布、代谢和排泄(ADME)特性。根据Lipinski五规则,口服药物的分子量应该小于500道尔顿。分子量过大会导致渗透性差,难以穿过生物膜；分子量过小可能无法形成足够的相互作用,导致活性不足。分子量的计算比较简单,就是把分子中所有原子的原子量加起来,RDKit和其他化学信息学库都提供了准确的计算功能。

氢键供体(Hydrogen Bond Donor)和氢键受体(Hydrogen Bond Acceptor)的数量反映了分子形成氢键的能力。氢键是药物-受体相互作用中的重要力量,对于分子的活性和选择性都有影响。一般来说,氢键供体的数量应该不超过5个,氢键受体的数量应该不超过10个(Lipinski规则)。氢键供体通常是指连接在电负性原子(如N、O)上的氢原子,氢键受体是指具有孤对电子的电负性原子。RDKit可以通过分析分子的结构来自动识别和计数氢键供体和受体。

可旋转键数量(Number of Rotatable Bonds)反映了分子的柔性。可旋转键越多,分子的构象空间越大,可能在结合时需要付出更多的熵代价。一般建议可旋转键数量不超过10个。可旋转键的定义通常是单键,不包括末端甲基的旋转,也不包括环内的键。RDKit提供了专门的函数来计算可旋转键的数量。

极性表面积(Topological Polar Surface Area, TPSA)是分子表面极性部分的面积之和,它与分子的渗透性和溶解性有关。TPSA的值越大,分子越亲水,越难穿过脂质膜。对于中枢神经系统药物,TPSA通常应该小于90平方埃；对于一般的口服药物,TPSA应该小于140平方埃。TPSA的计算基于原子类型和连接性,不需要三维结构,所以叫做"拓扑"极性表面积。

除了上述属性,还有很多其他的药物相关属性,比如合成可及性(Synthetic Accessibility, SA)、生物利用度(Bioavailability)、毒性(Toxicity)等。这些属性的预测方法各不相同,有的基于经验规则,有的基于统计模型,有的基于机器学习。在本研究中,我们主要关注QED属性,因为它是综合性的指标,而且RDKit提供了现成的计算函数,方便我们进行实验。当然,框架本身是通用的,可以很容易地扩展到其他属性的优化,只需要在属性预测模块中添加相应的计算函数即可。

在属性预测的实现上,我们使用了RDKit库提供的函数。RDKit是一个开源的化学信息学软件包\cite{landrum2013rdkit},它提供了丰富的分子处理和属性计算功能。我们的PropertyPredictor类封装了RDKit的调用,提供了一个统一的接口来预测各种属性。对于QED,我们使用\linebreak[0] rdkit.Chem.QED.qed()函数；对于LogP,我们使用\linebreak[0] rdkit.Chem.Crippen.\linebreak[0] MolLogP()函数；对于其他属性,也有对应的函数可用。这种方式的优点是准确可靠,因为这些函数都是经过广泛验证的；缺点是计算速度相对较慢,特别是对于大规模的数据集。在未来的工作中,可以考虑用机器学习模型来替代传统的计算方法,提高预测的速度和准确性。

\section{进化算法原理}

进化算法(Evolutionary Algorithm, EA)是一类受生物进化启发的优化算法,它们模拟了自然界中"优胜劣汰、适者生存"的进化过程,通过种群的迭代演化来寻找问题的最优解\cite{holland1992genetic}。进化算法的核心思想是维护一个由多个候选解组成的种群,通过选择、交叉、变异等操作,不断地产生新的候选解,并根据适应度函数来评估和筛选,最终收敛到高质量的解。进化算法的主要优势在于它们是全局搜索方法,不容易陷入局部最优；它们对问题的要求比较低,不需要梯度信息,可以处理不连续、非线性、甚至黑箱的优化问题；它们还具有很好的并行性,可以同时探索搜索空间的多个区域。

进化算法的基本流程包括以下几个步骤：初始化、评估、选择、遗传操作(交叉和变异)、更新、终止判断。首先是初始化阶段,随机生成一定数量的个体,构成初始种群。个体的表示方式取决于具体的问题,可以是二进制串、实数向量、排列、树结构等。在我们的分子优化问题中,个体就是分子的SMILES字符串。然后是评估阶段,对每个个体计算适应度值,适应度函数通常就是我们要优化的目标函数。在分子优化中,适应度可以是QED分数,或者其他药物属性的加权组合。

接下来是选择阶段,根据个体的适应度来选择哪些个体可以进入下一代,或者作为父代产生后代。选择的目的是施加选择压力,让优秀的个体有更多的机会传递它们的"基因"。常用的选择方法有轮盘赌选择(Roulette Wheel Selection)、锦标赛选择(Tournament Selection)、排名选择(Rank Selection)等。轮盘赌选择是按照适应度的比例来分配选择概率,适应度越高的个体被选中的概率越大；锦标赛选择是随机选取k个个体,从中选择适应度最高的那个；排名选择是先对个体按适应度排序,然后根据排名来分配选择概率。在我们的实现中,我们支持多种选择方法,可以通过参数来灵活配置。

遗传操作包括交叉(Crossover)和变异(Mutation)两个算子。交叉是把两个父代个体的部分基因组合起来,产生新的子代个体,它模拟了生物的有性繁殖过程。交叉的目的是结合不同父代的优点,产生更好的后代。变异的是对个体的某些基因进行随机的改变,它模拟了生物的基因突变现象。变异的作用是引入新的基因,增加种群的多样性,防止过早收敛。在传统的遗传算法中,交叉和变异都是基于预定义的规则,比如对于二进制串,可以使用单点交叉、多点交叉、均匀交叉等；对于实数向量,可以使用算术交叉、BLX-$\alpha$交叉等；变异可以是位翻转、高斯扰动等。

但是在我们的研究中,我们用大语言模型的生成操作替代了传统的变异算子。这是因为分子的SMILES表示是一种复杂的结构化字符串,传统的变异操作(如随机插入、删除、替换字符)很容易产生无效的SMILES,或者破坏分子的化学合理性。而大语言模型通过学习大量的化学知识,能够生成化学上有效且合理的分子结构,这是一个很大的优势。我们把当前分子的SMILES和优化目标用自然语言的形式提示给大语言模型,让它生成改进后的分子,这个过程既包含了变异的成分(产生新的结构),也包含了一定的导向性(朝着优化目标的方向)。

更新阶段是把新生成的子代个体加入到种群中,同时移除一些旧的个体,保持种群规模不变。更新的策略有多种,一种是世代更新(Generational Update),即完全用子代替换父代,只保留最优的个体(精英策略)；另一种是稳态更新(Steady-State Update),即每次只替换种群中最差的几个个体,种群的其他部分保持不变。世代更新的优点是进化速度快,缺点是可能丢失多样性；稳态更新的优点是保留了更多的历史信息,进化更加平稳,缺点是收敛速度可能较慢。在我们的实现中,我们默认使用稳态更新策略,因为它在实践中表现更好。

最后是终止判断,当满足某个终止条件时,算法停止运行。常见的终止条件包括达到最大的迭代代数、适应度达到阈值、种群多样性低于某个值、或者连续多代没有明显的改进等。在我们的实验中,我们设置了固定的最大代数作为终止条件,同时也记录了每代的最佳适应度和平均适应度,用来分析收敛的过程。

进化算法有很多变种,比如遗传算法(Genetic Algorithm, GA)、进化策略(Evolution Strategy, ES)、遗传编程(Genetic Programming, GP)、差分进化(Differential Evolution, DE)等。它们在个体的表示、遗传操作的设计、选择策略等方面有所不同,但是基本框架是相似的。我们的方法可以看作是一种特殊的遗传算法,其中变异算子被大语言模型替代,这是一种新颖的尝试,结合了符号人工智能(大语言模型)和进化计算的优势。

\section{大语言模型基础}

大语言模型(Large Language Model, LLM)是基于深度学习技术构建的自然语言处理模型,它们通过在大规模的文本语料库上进行预训练,学习到了语言的统计规律、语义知识、甚至是推理能力。大语言模型的核心架构是Transformer,它是由Vaswani等人在2017年提出的\cite{vaswani2017attention}。Transformer摒弃了传统的循环神经网络(RNN)和卷积神经网络(CNN)架构,完全基于自注意力机制(Self-Attention Mechanism)来建模序列中各个位置之间的依赖关系。这种设计使得模型能够并行处理整个序列,大大提高了训练的效率,同时也能够捕捉长距离的依赖关系,克服了RNN的梯度消失问题。

Transformer的基本组成单元是注意力头(Attention Head)、前馈神经网络(Feed-Forward Network)、层归一化(Layer Normalization)、以及残差连接(Residual Connection)。注意力机制的核心思想是为序列中的每个位置计算一个权重分布,表示该位置与其他位置的相关程度,然后根据这个权重分布来聚合信息。多头注意力(Multi-Head Attention)是使用多个注意力头并行计算,每个头关注不同的子空间,最后把结果拼接起来,这样能够增强模型的表达能力。Transformer有两种主要的变体：编码器-解码器架构(Encoder-Decoder),用于机器翻译、文本摘要等序列到序列的任务；纯解码器架构(Decoder-Only),用于文本生成、语言建模等任务。目前主流的大语言模型如GPT系列、LLaMA、Qwen等都是基于Decoder-Only架构的。

大语言模型的训练通常分为两个阶段：预训练(Pre-training)和微调(Fine-tuning)。预训练阶段是在海量的无标注文本数据上进行的,训练的目标是语言建模,即给定前面的词,预测下一个词的概率分布\cite{brown2020language}。通过这种方式,模型学习到了语言的语法结构、词汇的语义、以及一些世界知识。预训练的数据规模非常大,可以达到数千亿甚至数万亿的词元(Token),模型的参数量也从几亿到几千亿不等。预训练完成后,模型已经具备了很强的语言理解和生成能力,可以用于零样本(Zero-shot)或小样本(Few-shot)的任务。

微调阶段是在特定任务的标注数据上进行的,目的是让模型更好地适应具体的应用场景。微调的方法有多种,一种是全参数微调(Full Fine-tuning),即更新模型的所有参数,这种方法效果最好,但是计算成本很高；另一种是参数高效微调(Parameter-Efficient Fine-Tuning, PEFT),如LoRA(Low-Rank Adaptation)、Prefix Tuning、Adapter等,它们只更新少量的参数,大大降低了计算资源的需求。还有一种是提示微调(Prompt Tuning),即在输入中添加可学习的提示向量,引导模型完成特定的任务,这种方法不需要更新模型的参数,只需要优化提示向量即可。

在大语言模型的使用上,有两种主要的交互方式：一种是API调用,即通过云服务提供商的API接口来使用模型,如OpenAI的GPT-4 API、阿里云的通义千问API等；另一种是本地部署,即在本地服务器上运行开源的模型,如通过Ollama、Hugging Face Transformers等工具来加载和运行模型。API调用的优点是方便易用,不需要自己管理硬件和软件,缺点是数据隐私问题、网络延迟、以及费用成本。本地部署的优点是数据可控、响应速度快、没有使用限制,缺点是需要较强的硬件支持(特别是GPU显存),以及一定的技术门槛。

在本研究中,我们选择了本地部署的方式,使用Ollama工具来运行Qwen2.5-7B模型。Ollama是一个简单易用的大语言模型运行工具,它支持多种开源模型,提供了统一的API接口,并且自动处理模型的下载、量化、加载等细节。Qwen2.5是阿里巴巴通义实验室开发的开源大语言模型,7B表示它有70亿个参数,这是一个中等规模的模型,在消费级的GPU上就可以运行。虽然相比GPT-4这样的超大规模模型,Qwen2.5-7B的能力有限,但是对于我们的分子生成任务来说,它已经足够胜任,而且本地部署保证了实验的可重复性和数据的安全性。

大语言模型在生成文本时,会使用采样策略来决定下一个词的选择。最常用的采样方法是温度采样(Temperature Sampling),它通过调整softmax函数的温度参数来控制生成的随机性。温度值越高,概率分布越平坦,生成的结果越多样化,但是也越不可控；温度值越低,概率分布越尖锐,生成的结果越确定性高,但是可能缺乏创造性。除了温度采样,还有Top-k采样(只从概率最高的k个词中采样)、Top-p采样(核采样,从累积概率达到p的最小词集中采样)等方法。在我们的实验中,我们通过调整温度参数来控制大语言模型生成的多样性,并在第四章中进行了详细的分析和讨论。

大语言模型的一个重要特性是上下文学习(In-Context Learning),即通过在输入中提供示例,让模型在没有微调的情况下学会完成新的任务。这种现象也被称为少样本学习(Few-Shot Learning)。上下文学习的能力使得大语言模型非常灵活,用户可以通过设计合适的提示(Prompt)来引导模型完成各种任务,而不需要修改模型的参数。在我们的研究中,我们设计了多种提示策略,包括基础提示、带示例的提示、以及鼓励多样性的提示,这些都是利用了大语言模型的上下文学习能力。提示工程(Prompt Engineering)已经成为使用大语言模型的一项关键技术,好的提示设计可以显著提升模型的表现。

总的来说,大语言模型为我们提供了一种强大的工具,它能够理解和生成复杂的结构化文本,包括分子的SMILES表示。通过巧妙地设计提示和利用模型的生成能力,我们可以实现传统方法难以完成的分子优化任务。当然,大语言模型也有其局限性,比如可能会生成化学上无效的结构、对提示的敏感性较高、生成的随机性难以精确控制等,这些问题都需要在实际应用中加以注意和解决。

近年来,越来越多的研究开始探索大语言模型在化学和药物发现领域的应用。Galactica模型展示了在不进行专门微调的情况下完成分子性质预测等科学任务的能力\cite{taylor2022galactica}。ChemBERTa等专门的化学预训练模型也在分子属性预测任务上表现优异\cite{chithrananda2020chemberta}。这些工作为大语言模型在分子优化中的应用奠定了理论基础。
